\resourceheading{10}{Flexible Peer Roles and Participation Routes}{Social interaction, group work \& executive-function access\par}{Teachers, students, families, and activity leaders}{Kit for Kids lesson design, EDU486 camp \& Module 3 communication planning}

\vspace{-0.5cm}
\toolsection{What This Resource Is}

The same principle of distributed knowledge applies inside group work: belonging should not depend on one pace, posture, communication mode, sensory tolerance, or public role. This tool combines the Kit for Kids peer actions (\textit{invite, wait, ask, and respect}) with visible, changeable roles and explicit permission to observe, pass, pause, and return \citep{oar2026kit,oar2026tips}. The goal is to make the intellectual and social work of the group accessible while preserving choice, challenge, and meaningful contribution. Peer interaction becomes a rich communication context when students can initiate, direct, disagree, ask questions, repair misunderstandings, revise ideas, and influence shared work \citep{downing2015classroom}. Group systems therefore need to teach peers how to support without taking over while also requiring adults to redesign barriers such as noise, rushed timing, unclear roles, competition for materials, or an expectation that everyone respond aloud. Camp practice reinforced the same principle. Participation included measuring, observing, recording, speaking, partnering, stepping back, and returning rather than treating continuous visible engagement as the only acceptable route \citep{campEvidence2026}. A flexible role system makes those routes explicit before a learner has to struggle to create one.

\toolsection{Flexible Role Menu}

% \small
\renewcommand{\arraystretch}{1.0}
\setlength{\tabcolsep}{3pt}

\begin{longtable}{
>{\bfseries\RaggedRight\arraybackslash}p{0.90in}
>{\RaggedRight\arraybackslash}p{2.40in}
>{\RaggedRight\arraybackslash}p{3.60in}}

\tableheaderthree{Role}{Possible Contribution}{Access Routes and Safeguards}

Starter/\newline modeler &
Begin one step, demonstrate a possible move, or make an idea visible &
Use materials, a photograph, gesture, drawing, code, AAC, or partner support; another student may begin \\

\rowcolor{SoftGray}
Evidence finder &
Notice, collect, compare, locate, or point to information &
Use direct contact or non-contact observation, camera, magnifier, label, pointer, or partner handling; sensory contact is never required \\

Recorder/\newline mapper &
Preserve ideas, changes, questions, evidence, or group decisions &
Write, type, draw, photograph, arrange symbols, dictate, or direct a partner without requiring handwriting or continuous note taking \\

\rowcolor{SoftGray}
Connector/\newline questioner &
Ask how ideas fit, request clarification, invite another perspective, or identify what is still unknown &
Use speech, a prepared card, AAC, private note, pointing, writing, or a partner-read question \\

Materials/\newline access lead &
Organize tools and help notice whether everyone reaches the work &
Ask before moving another person's materials, seat, device, or communication system; the role supports access rather than controlling peers \\

\rowcolor{SoftGray}
Reviser/\newline debugger &
Test an idea, disagree, locate a problem, or propose another route &
``Not that,'' ``try again,'' ``I disagree,'' and ``something else'' remain valued contributions \\

Presenter/\newline demonstrator &
Share or demonstrate the group's learning for another audience &
Speak, show, point, play media, hold proof, use AAC, co-present, answer privately, or present another way \\

\rowcolor{SoftGray}
Observer/\newline returner &
Watch first, track the process, pause participation, and enter later through a named step &
The role includes a visible re-entry route and cannot become a holding place that excludes the learner from substantive intellectual work \\

\end{longtable}

\normalsize

\toolsection{Group Start Script}

Display shared goals, known materials, approximate time, and the role menu before the group starts. Say:

\begin{quote}
\textit{Choose a way to begin. You may take one role, combine roles, ask for a different one, observe first, or pass and return. Ask before helping. Wait for a response. Keep communication tools with their owner. A strong group changes its plan when someone identifies an access barrier.}
\end{quote}

Make useful group language available before students need it: \textit{my turn, your turn, wait, help, show me, not yet, stop, different role, I disagree, I have an idea, you misunderstood,} and \textit{I am ready to return}. These messages may be spoken, written, pointed to, signed, selected through AAC, or communicated through another reliable route. Use a visible turn sequence if it supports access; do not require students to contribute in only one way.

% \small
\toolsection{When Group Work Breaks Down}

\renewcommand{\arraystretch}{1.0}
\setlength{\tabcolsep}{3pt}

\begin{tabularx}{\textwidth}{
>{\RaggedRight\arraybackslash}p{1.35in}
>{\RaggedRight\arraybackslash}p{2.25in}
>{\RaggedRight\arraybackslash}X}

\rowcolor{PrimaryRed}
\color{white}What Adults Notice &
\color{white}Possible Barrier to Investigate &
\color{white}Support Match \\

One student controls tools or talk &
Roles or turn-taking are unclear; helping has become takeover &
Return tool ownership, offer role and turn choices, teach ask-and-wait \& create an independent route \\

\rowcolor{SoftGray}
Student withdraws, pauses, or leaves &
Noise, crowding, uncertainty, processing demand, conflict, fatigue, pain, or an unwanted role &
Offer space, a lower-load route, a private check-in, another group size or role, and a clear step for returning \\

Group rushes one member &
Speed is treated as competence and wait time is not protected &
Remove race language, offer flexible timing, allow asynchronous contribution \& model real waiting \\

\rowcolor{SoftGray}
An idea is ignored &
Communication mode, status, or group dynamics determine whose evidence counts &
Record proposals visibly, confirm meaning, provide private or anonymous routes, and require decisions to reference evidence, not popularity \\

Peer ``helps'' without consent &
Assistance has been praised without boundaries or the learner's control has been displaced &
Ask first, offer choices, accept ``no,'' return control of the task or tool, and repair the interaction when needed \\

\end{tabularx}

\normalsize
\toolsection{Review With Students}

Review the group process with the same attention given to the final product. Ask which roles were meaningful, whose ideas changed the work, whether anyone had to perform a preferred style to participate, where communication became difficult, and what adults should redesign. Record contribution through multiple forms. Constructing competence means designing environments that expect meaningful contribution while supplying the access necessary to make it possible \citep{jorgensen2018competence}. Puzzle Plan records whose contribution changed the product; EQUITAS asks how the role structure distributes influence and supports inclusion.

\reflectionbox{\textbf{Practice artifact, camp role trail:} Camp participation often changed form without losing intellectual value. During measurement, one youth self-selected the role of ``commentator'' and watched for visible litter. Identity Beads recognized contributions such as designer, environmental advocate, altruist, brave scientist, motivator, investigator, inventor, storyteller, or another youth-authored identity without ranking one above another. Policy work introduced mayor, parks, tourism, restaurant, and environmental-scientist perspectives while preserving vote. At the final showcase, youth could demonstrate, point, answer, ask a visitor a question, hold evidence, or partner with another speaker. Across these activities, roles carried real intellectual and social work and remained changeable.}

\adaptationbox{Peer structures cannot compensate for inaccessible instruction. Adults still have to address noise, pace, material layout, language demands, conflict, social status, and whose communication is recognized. Do not create a permanent helper, spokesperson, or observer role for a learner. Revisit roles across activities, protect private and bilingual communication, and treat refusal, withdrawal, or a request for individual work as information about access. A participation route is successful when it preserves both belonging and influence.}

% \sourcebox{\scriptsize OAR peer-education resources \citep{oar2026kit,oar2026tips}; rich communicative environments \citep{downing2015classroom}; communicative competence \citep{jorgensen2018competence}; de-identified camp evidence \citep{campEvidence2026}.}